% GENERATED FILE. Edit canonical lessons, then run npm run generate:books.
% canonical-source-hash: fnv1a64:2a81891cbae9f20b
% canonical-lessons: JA-C525-honmono, JA-C525-nisemono, JA-C525-marude, JA-C525-tatoeru, JA-C525-izure

\chapter{Real or Fake}
\label{ch:ja-real-or-fake}
\hlchaptermodality{525}

\begin{chapteropening}
\textbf{By the end of this chapter:} \emph{I can compare and say what would happen.}

Complete the last lesson of chapter 525: I can compare and say what would happen.
\end{chapteropening}

\section[\texorpdfstring{\ja{ほんもの} (honmono)}{ほんもの (honmono)}]{\ja{ほんもの} (honmono) — the real thing; genuine}

\label{lesson:JA-C525-honmono}



\begin{quote}
Before the new one: say the Japanese for reasoning, then the Japanese for to be realized.
\end{quote}

\subsection*{You'll want to know: \ja{ほんもの}}
\textbf{\ja{ほんもの}} — \emph{honmono} — \textquotedblleft{}the real thing; genuine\textquotedblright{}.

\begin{quote}
\textbf{\ja{もし} \ja{ほんものを} \ja{かったら、たかいでしょう。}} — \emph{moshi honmono o kattara, takai deshō} — \textquotedblleft{}If you buy the real thing, it will probably be expensive.\textquotedblright{}
\end{quote}

\begin{cousinweb}
\textbf{\ja{ほんもの}} (\emph{honmono}) mixes the two layers, native and Sino-Japanese: Sino-Japanese \emph{hon}, \textquotedblleft{}true\textquotedblright{} (Middle Chinese \emph{pwonX}), borrowed from Chinese and native \emph{mono}, \textquotedblleft{}thing\textquotedblright{} (JMdict, KANJIDIC2; Middle Chinese in Baxter's transcription).
\end{cousinweb}

\begin{grammarlens}[title={What you've built}]
One more word for saying what happens if.
\end{grammarlens}

\subsection*{Your turn}
\begin{itemize}
\raggedright
  \item \emph{Say it:} \emph{honmono}
  \item \emph{Say it:} \emph{honmono}, once more
  \item \emph{Say it:} \emph{moshi honmono o kattara, takai deshō} — If you buy the real thing, it will probably be expensive.
  \item \emph{Recall:} say the Japanese for reasoning, then the Japanese for to be realized, then say \emph{honmono} again
\end{itemize}

\begin{tcolorbox}[breakable,colback=teal!4,colframe=teal!35!black,title={Before you move on}]
What does \ja{ほんもの} mean? (\textquotedblleft{}The real thing; genuine\textquotedblright{}.) Say it once more.
\end{tcolorbox}
\section[\texorpdfstring{\ja{にせもの} (nisemono)}{にせもの (nisemono)}]{\ja{にせもの} (nisemono) — a fake; a forgery}

\label{lesson:JA-C525-nisemono}



\begin{quote}
Before the new one: say the Japanese for to be realized, then the Japanese for the real thing.
\end{quote}

\subsection*{You'll want to know: \ja{にせもの}}
\textbf{\ja{にせもの}} — \emph{nisemono} — \textquotedblleft{}a fake; a forgery\textquotedblright{}.

\begin{quote}
\textbf{\ja{もし} \ja{にせものが} \ja{みつかったら、けいさつに} \ja{しらせて} \ja{ください。}} — \emph{moshi nisemono ga mitsukattara, keisatsu ni shirasete kudasai} — \textquotedblleft{}If you find a fake, please inform the police.\textquotedblright{}
\end{quote}

\begin{cousinweb}
\textbf{\ja{にせもの}} (\emph{nisemono}) is a native Japanese compound (wago), inherited rather than borrowed from Chinese: built from \emph{nise}, \textquotedblleft{}false\textquotedblright{} and \emph{mono}, \textquotedblleft{}thing\textquotedblright{} (JMdict, KANJIDIC2).
\end{cousinweb}

\begin{grammarlens}[title={What you've built}]
One more word for saying what happens if.
\end{grammarlens}

\subsection*{Your turn}
\begin{itemize}
\raggedright
  \item \emph{Say it:} \emph{nisemono}
  \item \emph{Say it:} \emph{nisemono}, once more
  \item \emph{Say it:} \emph{moshi nisemono ga mitsukattara, keisatsu ni shirasete kudasai} — If you find a fake, please inform the police.
  \item \emph{Recall:} say the Japanese for to be realized, then the Japanese for the real thing, then say \emph{nisemono} again
\end{itemize}

\begin{tcolorbox}[breakable,colback=teal!4,colframe=teal!35!black,title={Before you move on}]
What does \ja{にせもの} mean? (\textquotedblleft{}A fake; a forgery\textquotedblright{}.) Say it once more.
\end{tcolorbox}
\section[\texorpdfstring{\ja{まるで} (marude)}{まるで (marude)}]{\ja{まるで} (marude) — (not) in the least; completely}

\label{lesson:JA-C525-marude}



\begin{quote}
Before the new one: say the Japanese for the real thing, then the Japanese for a fake.
\end{quote}

\subsection*{You'll want to know: \ja{まるで}}
\textbf{\ja{まるで}} — \emph{marude} — \textquotedblleft{}(not) in the least; completely\textquotedblright{}.

\begin{quote}
\textbf{\ja{もし} \ja{ふゆに} \ja{いったら、うみには} \ja{まるで} \ja{だれも} \ja{いません。}} — \emph{moshi fuyu ni ittara, umi ni wa marude daremo imasen} — \textquotedblleft{}If you go in winter, there is nobody at all at the sea.\textquotedblright{}
\end{quote}

\begin{cousinweb}
\textbf{\ja{まるで}} (\emph{marude}) is a native Japanese word (wago), inherited rather than borrowed from Chinese, and usually written in kana; its older origin is not settled (JMdict).
\end{cousinweb}

\begin{grammarlens}[title={What you've built}]
One more word for saying what happens if.
\end{grammarlens}

\subsection*{Your turn}
\begin{itemize}
\raggedright
  \item \emph{Say it:} \emph{marude}
  \item \emph{Say it:} \emph{marude}, once more
  \item \emph{Say it:} \emph{moshi fuyu ni ittara, umi ni wa marude daremo imasen} — If you go in winter, there is nobody at all at the sea.
  \item \emph{Recall:} say the Japanese for the real thing, then the Japanese for a fake, then say \emph{marude} again
\end{itemize}

\begin{tcolorbox}[breakable,colback=teal!4,colframe=teal!35!black,title={Before you move on}]
What does \ja{まるで} mean? (\textquotedblleft{}(not) in the least; completely\textquotedblright{}.) Say it once more.
\end{tcolorbox}
\section[\texorpdfstring{\ja{たとえる} (tatoeru)}{たとえる (tatoeru)}]{\ja{たとえる} (tatoeru) — to compare (to); to liken}

\label{lesson:JA-C525-tatoeru}



\begin{quote}
Before the new one: say the Japanese for a fake, then the Japanese for (not) in the least.
\end{quote}

\subsection*{You'll want to know: \ja{たとえる}}
\textbf{\ja{たとえる}} — \emph{tatoeru} — \textquotedblleft{}to compare (to); to liken\textquotedblright{}.

\begin{quote}
\textbf{\ja{もし} \ja{いのちを} \ja{はなに} \ja{たとえたら、どの} \ja{はなでしょう。}} — \emph{moshi inochi o hana ni tatoetara, dono hana deshō} — \textquotedblleft{}If you compared life to a flower, which flower would it be?\textquotedblright{}
\end{quote}

Polite \textbf{\ja{たとえます}}, \emph{tatoemasu}; an ichidan verb. Past (polite): \textbf{\ja{しじんが} \ja{こいを} \ja{はなに} \ja{たとえました。}}, \emph{shijin ga koi o hana ni tatoemashita}.

\begin{cousinweb}
\textbf{\ja{たとえる}} (\emph{tatoeru}) is a native Japanese word (wago), inherited rather than borrowed from Chinese; its kanji was attached later as a kun'yomi, and its older origin is not settled (JMdict, KANJIDIC2).
\end{cousinweb}

\begin{grammarlens}[title={What you've built}]
One more word for saying what happens if.
\end{grammarlens}

\subsection*{Your turn}
\begin{itemize}
\raggedright
  \item \emph{Say it:} \emph{tatoeru}
  \item \emph{Say it:} \emph{tatoeru}, once more
  \item \emph{Say it:} \emph{moshi inochi o hana ni tatoetara, dono hana deshō} — If you compared life to a flower, which flower would it be?
  \item \emph{Recall:} say the Japanese for a fake, then the Japanese for (not) in the least, then say \emph{tatoeru} again
\end{itemize}

\begin{tcolorbox}[breakable,colback=teal!4,colframe=teal!35!black,title={Before you move on}]
What does \ja{たとえる} mean? (\textquotedblleft{}To compare (to); to liken\textquotedblright{}.) Say it once more.
\end{tcolorbox}
\section[\texorpdfstring{\ja{いずれ} (izure)}{いずれ (izure)}]{\ja{いずれ} (izure) — sooner or later; some day}

\label{lesson:JA-C525-izure}



\begin{quote}
Before the new one: say the Japanese for (not) in the least, then the Japanese for to compare (to).
\end{quote}

\subsection*{You'll want to know: \ja{いずれ}}
\textbf{\ja{いずれ}} — \emph{izure} — \textquotedblleft{}sooner or later; some day\textquotedblright{}.

\begin{quote}
\textbf{\ja{いずれ} \ja{おとなに} \ja{なったら、むすこも} \ja{わかるでしょう。}} — \emph{izure otona ni nattara, musuko mo wakaru deshō} — \textquotedblleft{}Some day, when he grows up, my son will understand too.\textquotedblright{}
\end{quote}

\begin{cousinweb}
\textbf{\ja{いずれ}} (\emph{izure}) is a native Japanese word (wago), inherited rather than borrowed from Chinese, and usually written in kana; its older origin is not settled (JMdict).
\end{cousinweb}

\begin{grammarlens}[title={What you've built}]
One more word for saying what happens if.
\end{grammarlens}

\subsection*{Your turn}
\begin{itemize}
\raggedright
  \item \emph{Say it:} \emph{izure}
  \item \emph{Say it:} \emph{izure}, once more
  \item \emph{Say it:} \emph{izure otona ni nattara, musuko mo wakaru deshō} — Some day, when he grows up, my son will understand too.
  \item \emph{Recall:} say the Japanese for (not) in the least, then the Japanese for to compare (to), then say \emph{izure} again
\end{itemize}

\begin{tcolorbox}[breakable,colback=teal!4,colframe=teal!35!black,title={Before you move on}]
What does \ja{いずれ} mean? (\textquotedblleft{}Sooner or later; some day\textquotedblright{}.) Say it once more.
\end{tcolorbox}
